\clearpage
\section*{本章符号索引}
\markright{本章符号索引}

各符号在正文首次出现时解释。下表按已知输入、实际决策与派生变量分类；时间统一换算为秒，功率、能量与温差分别采用W、J和K。

\begin{longtable}{@{}p{5.3cm}p{6.0cm}>{\raggedright\arraybackslash}p{2.5cm}@{}}
\caption{第三章主要符号}\label{tab:symbols}\\
\toprule
符号 & 含义 & 量纲／类别 \\
\midrule
\endfirsthead
\toprule
符号 & 含义 & 量纲／类别 \\
\midrule
\endhead
$\mathcal A,G_a,V_a,E_a$ & 应用集合、各应用的Stage依赖图 & 输入集合 \\
$\mathcal C,\mathcal R_c,\mathcal I_c,\mathcal E_c$ & 场景、请求、Stage实例及依赖 & 输入集合 \\
$\mathcal J,\tau(c,i),r(c,i),\mathcal O_{cr}$ & 固定Stage身份、实例归属、必需输出 & 输入映射 \\
$r_{cr},d^{\rm rel}_{cr},D_{cr},W_c$ & 释放、期限与固定评价窗口 & s \\
$p_c$ & 完整场景的统计权重 & 无量纲 \\
$\mathcal U_C,\mathcal U_G,\mathcal U_N$ & CPU、GPU与NPU执行单元 & 输入集合 \\
$\mathcal H,\gamma(u),\mathcal K_h,L_h$ & 频率域、所属域、工作点、最低档 & 输入集合 \\
$f_h(k),v_h(k)$ & 配套频率与电压 & Hz、V \\
$\mathcal M_{ci},u(m),\kappa(m)$ & 合法模式、具体设备及工作点 & 输入映射 \\
$d^{\rm ref}_{\tau u},f^{\rm ref}_u,\alpha_{\tau u}$ & 参考时长、频率和响应参数 & s、Hz、比例 \\
$d_\tau(m),q_\tau(m)$ & 模式持续时间、活动功率 & s、W \\
$\mathcal B,\pi(u),\xi_{cg}(t)$ & 功率计量组、单元归属与控制状态 & 输入／派生 \\
$B_g,\lambda_{gn},T_g^{\rm ref}$ & 基准、温度漏电系数和参考温度 & W、W/K、$^\circ$C \\
$\mathcal N,C,G,T_a,T_c^0,T_n^{\max}$ & 节点、热容、导热、初温与温限 & J/K、W/K、$^\circ$C \\
$P_{{\rm ext},c},E_{{\rm ext},c}$ & 网络外公共功率及能量 & W、J \\
\midrule
$u_\tau,m_{ci},s_{ci}$ & 共享设备部署、模式、起点 & 决策 \\
$F_{ci},I_{ci},\iota_{ci},k_{ch}(t)$ & 结束、执行区间、活动与域状态 & 派生 \\
$P_{cg},p^{\rm heat}_{cn},\theta_c,E_c,\bar E$ & 组功率、节点注热、温升和能量 & W、K、J \\
\midrule
$\mathcal G_{cim},\mathcal Q_{ci},\mathcal Q_c$ & 有限起点与Stage选项 & 输入／派生 \\
$t_o,f_o,d_o,q_o,u_o,h_o,\kappa_o$ & 选项执行参数与控制要求 & 已知系数 \\
$e_{c\ell},\delta_{c\ell},\chi_{co\ell}$ & 固定事件、片长及活动覆盖 & s、0/1 \\
$x_{\tau u},y_{co}$ & 共享部署、选项选择 & 二元变量 \\
$\rho_{ch\ell k},w_{c\ell\omega}$ & 控制域和联合状态选择 & 二元变量 \\
$\Omega,v_{c\ell}$ & 相容联合状态、活动选项向量 & 派生 \\
$\theta_{c\ell},\Delta E_{c\ell}$ & 片段端点温升与段能量 & K、J \\
$A_\omega,b_\omega,D_c,\beta_\omega,\ell_\omega$ & 热方程及电功率系数 & 派生系数 \\
$\Phi,\psi,\Gamma,\mu,\zeta,\nu$ & 状态传播与积分参数 & 派生系数 \\
$H_c,\overline P_c,M^T,M^E$ & 数值条件等式所用有效界 & 派生界 \\
$\mathcal H^{(j)},U$ & 热必要行集合、安全方案能量上界 & 求解记录 \\
\bottomrule
\end{longtable}
